% !TeX program = xelatex
% Standalone source: compile twice with XeLaTeX. No external figure files.
% Self-contained Letter; numerical results validated against the stored calculations.
% Figure coordinates below are empirical CDFs of the stored 200,000 series.
% Abstract, introduction and conclusion revised to clarify scope and motivation.
% Revision 2026-11-06: reading guidance is maintained in the repository README.
\documentclass[10pt,a4paper]{article}
\usepackage[a4paper,left=21mm,right=21mm,top=19mm,bottom=18mm,headheight=13pt,headsep=5mm,footskip=8mm]{geometry}
\usepackage{fontspec}
\setmainfont{Noto Serif CJK JP}[ItalicFont={Noto Serif CJK JP},ItalicFeatures={FakeSlant=0.15}]
\setsansfont{Noto Sans CJK JP}
\XeTeXlinebreaklocale "ja"
\XeTeXlinebreakskip=0pt plus 0.1em
\usepackage{amsmath,amssymb,booktabs,array}
\usepackage{xcolor,graphicx,pgfplots}
\pgfplotsset{compat=1.18}
\usepackage{titlesec,fancyhdr,caption,hyperref}
\definecolor{accent}{HTML}{176B80}
\definecolor{compare}{HTML}{B34D31}
\definecolor{muted}{HTML}{555555}
\hypersetup{unicode=true,colorlinks=true,linkcolor=accent,urlcolor=accent,citecolor=accent,
 pdftitle={記録から見えない状態を推定する：Letter版},pdfauthor={},
 pdfsubject={地下チンチロから潜在状態の推定と実社会への応用へ}}
\urlstyle{same}
\linespread{1.08}
\setlength{\parindent}{1em}
\setlength{\parskip}{2.5pt}
\setlength{\abovedisplayskip}{6pt plus 1pt minus 1pt}
\setlength{\belowdisplayskip}{6pt plus 1pt minus 1pt}
\setlength{\abovedisplayshortskip}{3pt}
\setlength{\belowdisplayshortskip}{4pt}
\renewcommand{\arraystretch}{1.13}
\titleformat{\section}{\sffamily\bfseries\normalsize\color{accent}}{\thesection}{0.65em}{}
\titlespacing*{\section}{0pt}{8pt}{4pt}
\titleformat{\subsection}{\sffamily\bfseries\small}{\thesubsection}{0.6em}{}
\titlespacing*{\subsection}{0pt}{6pt}{3pt}
\captionsetup{font=small,labelfont=bf,labelsep=quad,skip=3pt,hypcap=false}
\renewcommand{\figurename}{図}
\renewcommand{\tablename}{表}
\renewcommand{\refname}{参考文献}
\pagestyle{fancy}
\fancyhf{}
\fancyhead[L]{\sffamily\scriptsize 記録から見えない状態を推定する}
\fancyhead[R]{\sffamily\scriptsize Letter版\enspace |\enspace 2026年11月6日}
\fancyfoot[C]{\sffamily\small\thepage}
\renewcommand{\headrulewidth}{0.3pt}
\newcommand{\given}{\mid}
\newcommand{\BF}{\mathrm{BF}}
\newcommand{\F}{\mathcal F}
\newcommand{\Hi}{\mathrm{Hi}}
\newcommand{\Lo}{\mathrm{Lo}}

\begin{document}
\thispagestyle{fancy}
\begin{center}
{\sffamily\bfseries\fontsize{19}{25}\selectfont 記録から見えない状態を推定する\par}
\vspace{2mm}
{\sffamily\large 地下チンチロから潜在状態の推定と実社会への応用へ\par}
\end{center}
\vspace{-2mm}
\noindent\textbf{要旨}\enspace
直接観測できない状態や指標を、観測可能な記録と背景情報からどこまで数値的に推定できるかを検討する。
『カイジ』の地下チンチロの二項モデルに基づくベイズ推定を土台に、局面などの外的要因、観測間の依存関係、潜在状態の変化を取り込む推論へ考察を拡張した。
固定戦略の判定では、事前確率0.100、判定閾値0.990の条件で、期間無制限の誤警報上界$1/891$（約0.112\%）を得た。
理解度推定の応用例であるベイズ知識追跡（Bayesian Knowledge Tracing; BKT）では、同じ8回答への対数尤度差が約$4.432\times10^{-5}$の2候補でも、習得確率は0.984と0.369に分かれた。
潜在状態の推定と観測予測の保証を区別し、履歴・出題条件を含む条件付き予測の逐次監視を構成した。
さらに不確かさと因果効果の評価条件を整理し、大規模言語モデル（LLM）の品質管理など、直接観測しにくい対象を扱う分野への応用と実証の見通しを示す。

\section{地下チンチロを足場に、見えない状態の推定へ}
本稿は、直接観測できない状態や指標を観測可能な情報から推定する方法と、その実社会への応用可能性を検討する。
出発点とするのは、『カイジ』の地下チンチロ\cite{kaiji}を題材とした、二項モデルに基づくベイズ推定である。
出目という具体的な記録から、直接には分からない賽の使用戦略を推定する題材は、読者になじみのある導入であると同時に、生成規則を明示して推論を検証する基盤となる。

単一賽の二値観測から始め、局面に応じた使用率、3個一式で共有される種類、振り直しの条件へモデルを広げると、同じ見かけの結果でも証拠の強さは変わる。
ここで推定するのは、どの戦略仮説が記録を説明するかという確率であり、条件ごとの尤度から更新量を定める。
物語における三好の記録とカイジの着眼点を、情報を保存する役割と、比較すべき説明を構成する役割として捉える。

この推論を他の対象へ移す例として、学習時の理解度推定を扱う。
正誤だけでは分かりにくい状態を、出題条件や、自信度・解答時間などの情報と合わせてどう推定するかが実務上の動機にある。
本稿ではまずBKTの基本形を用い、観測への当てはまり、潜在状態の推定、次回正答の予測を区別して検証する。
自信度等を追加した際の精度向上は、今後の検証課題とする。

さらに、観測条件と結果の関連を利用する推定から、条件を変える介入の因果効果を検討する段階までを射程に含める。
以下では、モデルの下で計算・検証した結果と、理解度推定やLLM品質管理への応用・実証計画を分けて述べる。

\section{観測単位と共有状態が尤度を変える}
通常賽では1〜6、456賽では4〜6がそれぞれ等確率と仮定する。
単一賽の出目を$\Hi=\{4,5,6\}$、$\Lo=\{1,2,3\}$へまとめる。
$H_0$を常に通常賽、$H_1$を各観測で確率$\eta$により456賽を選ぶ戦略とすると、
\begin{equation}
 P(\Hi\given H_0)=\tfrac12,\qquad
 P(\Hi\given H_1)=\tfrac{1+\eta}{2}.
\end{equation}
履歴$D_{1:t}$に対する仮説確率を$\pi_t=P(H_1\given D_{1:t})$、オッズを$O_t=\pi_t/(1-\pi_t)$とおけば、
\begin{equation}
 O_t=O_{t-1}B_t,\qquad
 B_t=\frac{P(D_t\given H_1,D_{1:t-1})}{P(D_t\given H_0,D_{1:t-1})},\qquad
 \pi_t=\frac{O_t}{1+O_t}.
 \label{eq:odds}
\end{equation}
例えば$\eta=0.900$では、$\Hi$はオッズを1.900倍、$\Lo$は0.100倍にする。
反対向きの証拠は対称ではなく、確率への一定量の加減算にもならない。

実際の一式3個で同じ種類を共有し、手番内の振り直し中も維持する場合は、全投擲履歴$T_t$の尤度を
\begin{equation}
 P(T_t\given H_1,C_t)
 =\eta(C_t)P_{456}(T_t)+\{1-\eta(C_t)\}P_F(T_t)
 \label{eq:shared}
\end{equation}
とする。$C_t$は結果を見る前に分かる局面、$P_F,P_{456}$は種類を固定した履歴分布である。
共通の種類を周辺化した後の出目は一般に独立ではない。
例えば$\eta=0.900$で1投の全個体が$\Hi$となる確率は0.913であり、各賽の周辺確率を掛けた$0.950^3\simeq0.857$とは異なる。
また、通常時の使用率を$\eta(C_t)=0$とすれば、そこでの結果は$B_t=1$となる。
通常時の負けが疑いを下げるかどうかも、比較する戦略仮説によって変わる。

この例での精密化とは、記録を細分化することだけではない。
どの観測が同じ原因を共有し、いつ潜在状態が切り替わるかを定め、モデルに合った尤度を使うことである。
条件と結果の関連を取り込む段階と、条件を変えた場合の因果効果を求める段階も分けて考える。

\newpage
\section{逐次判定：確信度と誤警報率を分ける}
各手番で一式を独立に選ぶ一定使用率のモデルでは、第1投が全個体$\Hi$かを$X_t\in\{0,1\}$で表すと、
\begin{equation}
 P_0(X_t=1)=\tfrac18,\quad P_\eta(X_t=1)=\tfrac{1+7\eta}{8},\quad
 B_t=\begin{cases}1+7\eta,&X_t=1,\\1-\eta,&X_t=0.\end{cases}
\end{equation}
この2仮説の比較では、$X_t$が全投擲履歴の尤度比を保持する。
累積ベイズ因子$M_n=\prod_{t=1}^nB_t$は$H_0$の下で初期値1の非負マルチンゲールである。
事前確率$\pi_0$、事後確率の判定閾値$h$を事前に固定すると、Villeの不等式から\cite{howard}
\begin{equation}
 A=\frac{h}{1-h}\frac{1-\pi_0}{\pi_0},\qquad
 P_0\!\left(\exists n:M_n\ge A\right)\le A^{-1}.
 \label{eq:ville}
\end{equation}
$\pi_0=0.100,h=0.990$なら$A=891$で、期間無制限の誤警報上界は$1/891$（約0.112\%）となる。
事後確率99.000\%を、そのまま誤警報率1.000\%と読むことはできない。
未到達経路を整数再帰で集計した有限期間の値を表\ref{tab:dice}に示す。

\begin{center}
\captionof{table}{手番ごとの判定性能（真・想定使用率ともに$\eta=0.900$）}\label{tab:dice}
\small
\begin{tabular}{rrr}
\toprule 上限手番数 & 通常賽のみ：誤警報率 & 混合戦略：検知率\\
\midrule 4 & 0.024\% & 69.332\%\\
10 & 0.041\% & 99.356\%\\
20 & 0.041\% & 99.994\%\\
\bottomrule
\end{tabular}
\end{center}
観測前に決まる局面に応じて$\eta$を変えても、条件付き期待倍率が1なら保証を保てる。
結果を見てから都合のよい局面・仮説を選び直す操作は、この保証に含まれない。
未知の$\eta$を扱う場合は、事前に固定した分布で尤度比を混合する方法がある。


\section{BKT：習得状態と正答予測を分ける}
BKT（Bayesian Knowledge Tracing）は、特定スキルの潜在習得状態$L_t\in\{0,1\}$を正誤$Y_t$から推定する\cite{corbett}。
固定戦略の仮説確率と異なり、状態は学習により変わる。
解答前を$q_t=P(L_t=1\given Y_{1:t-1})$、解答後を$p_t=P(L_t=1\given Y_{1:t})$とする。
未習得時の正答率を$G$、習得時の誤答率を$S$、回答後の学習遷移確率を$T$とすれば、忘却なしの基本形は
\begin{align}
 p_t&=\begin{cases}
 \dfrac{(1-S)q_t}{(1-S)q_t+G(1-q_t)},&Y_t=1,\\[3pt]
 \dfrac{Sq_t}{Sq_t+(1-G)(1-q_t)},&Y_t=0,
 \end{cases}\\[-2pt]
 q_{t+1}&=p_t+(1-p_t)T,\qquad
 r_{t+1}=P(Y_{t+1}=1\given Y_{1:t})=(1-S)q_{t+1}+G(1-q_{t+1}).
 \label{eq:bkt}
\end{align}
状態確率$p_t$と次回正答予測$r_{t+1}$は異なる量である。
また、$T$はモデルの遷移パラメータであり、特定の解説の因果効果を自動的に表すものではない。

\begin{center}
\captionof{table}{同じ8回答$[1,1,0,1,0,1,1,1]$に対する固定候補の比較}\label{tab:bkt}
\small\setlength{\tabcolsep}{5pt}
\begin{tabular*}{0.97\linewidth}{@{\extracolsep{\fill}}lrrrrrrr@{}}
\toprule 候補 & $G$ & $S$ & $T$ & $q_1$ & 対数尤度 & $p_8$ & $r_9$\\
\midrule 基準 & 0.250 & 0.100 & 0.150 & 0.200 & $-6.135$ & 0.984 & 0.891\\
別解 & 0.415 & 0.040 & 0.020 & 0.100 & $-6.135$ & 0.369 & 0.623\\
\bottomrule
\end{tabular*}
\end{center}
表\ref{tab:bkt}の対数尤度差は約$4.432\times10^{-5}$だが、状態推定は大きく分かれる。
分数による再帰と全潜在経路の列挙を照合しており、この差は丸め・近似計算に由来しない。
示したのは有限の1系列での曖昧さであり\cite{beck}、全観測列で分布が一致する構造的識別不能性の証明ではない。
実際、両者の次回予測は異なる。


\newpage
\section{将来回答で候補を比べ、観測予測を監視する}
\subsection{将来の追加情報は、何を区別できるか}
表\ref{tab:bkt}の2候補を固定し、同じスキル・忘却なし・毎回答後の学習継続という条件で、過去8回答から先を比較する。
各候補の次回予測を追加履歴で更新し、その予測尤度比を累積する。
過去8回答は候補の構成に使ったため、将来部分の尤度比は1から開始する。
各生成モデル20万系列、上限400問の独立なモンテカルロ計算により、正しい方向の尤度比が10倍へ初めて達するまでを評価した（図\ref{fig:future}）。
両フィルタは過去8回答後の各候補自身の$q_9$（基準0.986、別解0.382）から再開し、生成側の条件付き正答予測$r_t$から回答を直接生成した（潜在状態は周辺化）。
乱数はPCG64、種\texttt{2026110302}から生成側ごとに分岐し、分位点は経験分布の逆関数で定義した。

\begin{center}
\begin{tikzpicture}
\begin{axis}[
 width=\linewidth,height=49mm,scale only axis=false,
 xmin=0,xmax=400,ymin=0,ymax=100,
 xtick={0,100,200,300,400},ytick={0,25,50,75,100},
 xlabel={追加問題数},ylabel={10倍への累積到達率（\%）},
 tick label style={font=\scriptsize},label style={font=\small},
 axis line style={black!55},tick style={black!45},
 grid=major,grid style={black!10},
 legend style={font=\scriptsize,draw=none,fill=none,at={(0.98,0.08)},anchor=south east},
 legend cell align=left,clip=true]
% BEGIN_EMPIRICAL_CDF
\addplot[color=accent,line width=1pt,const plot,no marks] coordinates {
(0,0.00000000)
(1,0.00000000)
(2,0.00000000)
(3,0.00000000)
(4,0.00000000)
(5,0.00000000)
(6,0.00000000)
(7,0.00000000)
(8,0.00000000)
(9,0.00000000)
(10,0.00000000)
(11,0.00000000)
(12,0.00000000)
(13,2.28900000)
(14,5.86700000)
(15,7.99750000)
(16,10.86550000)
(17,13.79150000)
(18,16.56950000)
(19,19.34250000)
(20,22.13550000)
(21,24.82200000)
(22,27.36200000)
(23,29.95250000)
(24,32.11100000)
(25,32.72950000)
(26,33.61450000)
(27,34.64050000)
(28,35.79200000)
(29,37.08000000)
(30,38.50700000)
(31,39.85100000)
(32,41.27950000)
(33,42.77150000)
(34,44.36950000)
(35,45.95650000)
(36,47.50000000)
(37,49.07300000)
(38,50.64600000)
(39,52.10350000)
(40,52.42350000)
(41,52.87950000)
(42,53.45950000)
(43,54.11150000)
(44,54.90050000)
(45,55.73750000)
(46,56.59650000)
(47,57.46750000)
(48,58.40750000)
(49,59.42700000)
(50,60.40950000)
(51,61.39500000)
(52,62.40450000)
(53,63.40500000)
(54,64.35800000)
(55,64.56550000)
(56,64.84600000)
(57,65.19600000)
(58,65.62300000)
(59,66.12850000)
(60,66.68950000)
(61,67.30500000)
(62,67.88400000)
(63,68.49500000)
(64,69.19050000)
(65,69.91250000)
(66,70.61300000)
(67,71.36250000)
(68,72.08350000)
(69,72.74500000)
(70,72.87950000)
(71,73.06850000)
(72,73.31350000)
(73,73.62900000)
(74,73.97900000)
(75,74.39500000)
(76,74.82700000)
(77,75.27150000)
(78,75.72400000)
(79,76.19150000)
(80,76.67700000)
(81,77.17550000)
(82,77.70300000)
(83,78.23550000)
(84,78.74850000)
(85,79.19300000)
(86,79.28500000)
(87,79.44050000)
(88,79.63200000)
(89,79.85050000)
(90,80.12300000)
(91,80.41500000)
(92,80.68450000)
(93,80.98950000)
(94,81.33850000)
(95,81.67400000)
(96,82.03550000)
(97,82.40400000)
(98,82.77900000)
(99,83.15950000)
(100,83.51800000)
(101,83.58350000)
(102,83.70100000)
(103,83.83950000)
(104,84.00400000)
(105,84.19650000)
(106,84.42200000)
(107,84.61750000)
(108,84.86100000)
(109,85.10500000)
(110,85.37100000)
(111,85.65900000)
(112,85.90800000)
(113,86.18600000)
(114,86.47400000)
(115,86.75700000)
(116,86.81350000)
(117,86.89300000)
(118,87.00700000)
(119,87.13200000)
(120,87.28000000)
(121,87.43200000)
(122,87.59450000)
(123,87.79450000)
(124,87.98700000)
(125,88.20750000)
(126,88.41550000)
(127,88.62500000)
(128,88.84500000)
(129,89.08200000)
(130,89.27450000)
(131,89.31550000)
(132,89.37950000)
(133,89.46800000)
(134,89.56300000)
(135,89.67500000)
(136,89.79900000)
(137,89.93400000)
(138,90.07250000)
(139,90.21400000)
(140,90.36900000)
(141,90.53150000)
(142,90.68550000)
(143,90.84600000)
(144,91.02750000)
(145,91.19800000)
(146,91.23100000)
(147,91.28500000)
(148,91.34700000)
(149,91.42000000)
(150,91.51250000)
(151,91.60900000)
(152,91.72200000)
(153,91.83550000)
(154,91.95350000)
(155,92.08650000)
(156,92.21600000)
(157,92.34200000)
(158,92.48900000)
(159,92.66150000)
(160,92.81150000)
(161,92.93850000)
(162,92.97000000)
(163,93.01100000)
(164,93.05800000)
(165,93.11800000)
(166,93.19150000)
(167,93.27400000)
(168,93.35600000)
(169,93.45050000)
(170,93.54750000)
(171,93.64200000)
(172,93.73550000)
(173,93.83850000)
(174,93.94400000)
(175,94.06250000)
(176,94.16300000)
(177,94.18250000)
(178,94.22400000)
(179,94.27350000)
(180,94.32500000)
(181,94.37800000)
(182,94.44200000)
(183,94.50050000)
(184,94.58750000)
(185,94.65650000)
(186,94.73500000)
(187,94.82350000)
(188,94.90600000)
(189,94.99200000)
(190,95.08150000)
(191,95.17100000)
(192,95.18900000)
(193,95.21500000)
(194,95.24650000)
(195,95.28600000)
(196,95.33700000)
(197,95.39900000)
(198,95.45850000)
(199,95.51250000)
(200,95.57500000)
(201,95.64050000)
(202,95.71150000)
(203,95.78150000)
(204,95.85250000)
(205,95.93050000)
(206,96.00450000)
(207,96.01950000)
(208,96.03650000)
(209,96.06200000)
(210,96.09350000)
(211,96.13350000)
(212,96.17100000)
(213,96.21700000)
(214,96.27100000)
(215,96.32200000)
(216,96.36550000)
(217,96.41800000)
(218,96.47400000)
(219,96.52400000)
(220,96.58300000)
(221,96.64400000)
(222,96.65600000)
(223,96.67550000)
(224,96.69550000)
(225,96.72550000)
(226,96.75400000)
(227,96.80050000)
(228,96.83800000)
(229,96.87950000)
(230,96.92100000)
(231,96.96250000)
(232,97.01550000)
(233,97.06250000)
(234,97.11850000)
(235,97.16450000)
(236,97.20400000)
(237,97.24350000)
(238,97.25350000)
(239,97.27050000)
(240,97.28850000)
(241,97.30800000)
(242,97.33000000)
(243,97.36000000)
(244,97.39350000)
(245,97.42300000)
(246,97.45450000)
(247,97.49100000)
(248,97.52500000)
(249,97.56200000)
(250,97.60200000)
(251,97.64350000)
(252,97.68100000)
(253,97.69000000)
(254,97.70150000)
(255,97.72050000)
(256,97.73550000)
(257,97.76100000)
(258,97.78700000)
(259,97.81000000)
(260,97.82900000)
(261,97.85850000)
(262,97.88600000)
(263,97.91700000)
(264,97.94450000)
(265,97.97850000)
(266,98.01300000)
(267,98.03550000)
(268,98.03950000)
(269,98.04800000)
(270,98.05650000)
(271,98.06950000)
(272,98.09150000)
(273,98.10500000)
(274,98.12300000)
(275,98.14150000)
(276,98.16650000)
(277,98.18900000)
(278,98.21100000)
(279,98.23800000)
(280,98.26400000)
(281,98.29550000)
(282,98.32450000)
(283,98.32750000)
(284,98.33350000)
(285,98.34800000)
(286,98.35950000)
(287,98.37250000)
(288,98.38500000)
(289,98.39950000)
(290,98.42150000)
(291,98.43650000)
(292,98.45650000)
(293,98.48450000)
(294,98.50700000)
(295,98.52900000)
(296,98.54950000)
(297,98.56500000)
(298,98.56900000)
(299,98.57400000)
(300,98.58100000)
(301,98.59150000)
(302,98.60450000)
(303,98.61650000)
(304,98.63350000)
(305,98.64650000)
(306,98.66500000)
(307,98.68150000)
(308,98.70050000)
(309,98.71800000)
(310,98.73650000)
(311,98.75650000)
(312,98.78050000)
(313,98.79700000)
(314,98.80200000)
(315,98.80800000)
(316,98.81600000)
(317,98.82550000)
(318,98.84000000)
(319,98.85100000)
(320,98.85950000)
(321,98.86950000)
(322,98.88500000)
(323,98.90500000)
(324,98.92150000)
(325,98.93650000)
(326,98.95150000)
(327,98.96300000)
(328,98.97550000)
(329,98.97900000)
(330,98.98550000)
(331,98.99150000)
(332,98.99650000)
(333,99.00450000)
(334,99.01400000)
(335,99.02200000)
(336,99.02800000)
(337,99.04400000)
(338,99.05450000)
(339,99.06550000)
(340,99.07700000)
(341,99.08700000)
(342,99.10100000)
(343,99.11600000)
(344,99.11850000)
(345,99.12100000)
(346,99.12600000)
(347,99.13000000)
(348,99.13900000)
(349,99.14650000)
(350,99.15150000)
(351,99.16150000)
(352,99.17100000)
(353,99.17800000)
(354,99.18900000)
(355,99.20000000)
(356,99.20950000)
(357,99.21800000)
(358,99.23400000)
(359,99.23700000)
(360,99.23950000)
(361,99.24400000)
(362,99.25100000)
(363,99.25650000)
(364,99.26300000)
(365,99.27050000)
(366,99.27950000)
(367,99.28600000)
(368,99.29400000)
(369,99.30350000)
(370,99.31150000)
(371,99.32050000)
(372,99.33300000)
(373,99.34050000)
(374,99.34150000)
(375,99.34400000)
(376,99.34700000)
(377,99.35250000)
(378,99.35550000)
(379,99.35800000)
(380,99.36700000)
(381,99.37950000)
(382,99.38400000)
(383,99.39150000)
(384,99.39850000)
(385,99.40350000)
(386,99.40850000)
(387,99.41850000)
(388,99.42550000)
(389,99.43150000)
(390,99.43400000)
(391,99.43450000)
(392,99.43600000)
(393,99.44000000)
(394,99.44550000)
(395,99.45300000)
(396,99.45800000)
(397,99.46750000)
(398,99.47450000)
(399,99.47900000)
(400,99.48750000)
};
\addlegendentry{基準が真}
\addplot[color=compare,line width=1pt,const plot,no marks] coordinates {
(0,0.00000000)
(1,0.00000000)
(2,20.77650000)
(3,20.77650000)
(4,30.47700000)
(5,40.34050000)
(6,47.67800000)
(7,47.67800000)
(8,50.30650000)
(9,52.88200000)
(10,52.88200000)
(11,54.00450000)
(12,55.11850000)
(13,55.11850000)
(14,55.60650000)
(15,56.15300000)
(16,56.15300000)
(17,56.41300000)
(18,56.68700000)
(19,56.68700000)
(20,56.81500000)
(21,56.95000000)
(22,57.05300000)
(23,57.05300000)
(24,57.10950000)
(25,57.16050000)
(26,57.16050000)
(27,57.17650000)
(28,57.20900000)
(29,57.20900000)
(30,57.22200000)
(31,57.23850000)
(32,57.23850000)
(33,57.24650000)
(34,57.25700000)
(35,57.25700000)
(36,57.26150000)
(37,57.26400000)
(38,57.26800000)
(39,57.26800000)
(40,57.27000000)
(41,57.27150000)
(42,57.27150000)
(43,57.27250000)
(44,57.27300000)
(45,57.27300000)
(46,57.27350000)
(47,57.27400000)
(48,57.27400000)
(49,57.27450000)
(50,57.27450000)
(51,62.14850000)
(52,62.14850000)
(53,62.18850000)
(54,62.23950000)
(55,62.23950000)
(56,62.23950000)
(57,62.51100000)
(58,62.51100000)
(59,62.51700000)
(60,62.51700000)
(61,62.72300000)
(62,62.74400000)
(63,62.77050000)
(64,62.91950000)
(65,63.16550000)
(66,68.21800000)
(67,68.22600000)
(68,68.27200000)
(69,68.34600000)
(70,68.35900000)
(71,68.37150000)
(72,68.70400000)
(73,68.71450000)
(74,68.72550000)
(75,68.73450000)
(76,68.78950000)
(77,68.97900000)
(78,69.03250000)
(79,69.21400000)
(80,69.41650000)
(81,73.89900000)
(82,73.91450000)
(83,73.96450000)
(84,74.02350000)
(85,74.04100000)
(86,74.07900000)
(87,74.37400000)
(88,74.38750000)
(89,74.39650000)
(90,74.42200000)
(91,74.47850000)
(92,74.69650000)
(93,74.73050000)
(94,74.96050000)
(95,75.14350000)
(96,78.79450000)
(97,78.81800000)
(98,78.85850000)
(99,78.89950000)
(100,78.91450000)
(101,78.94000000)
(102,78.97150000)
(103,79.19200000)
(104,79.20150000)
(105,79.21850000)
(106,79.26650000)
(107,79.45100000)
(108,79.48200000)
(109,79.65800000)
(110,79.80750000)
(111,82.64400000)
(112,82.66300000)
(113,82.69150000)
(114,82.72150000)
(115,82.74800000)
(116,82.76800000)
(117,82.79400000)
(118,82.98000000)
(119,82.98750000)
(120,83.00850000)
(121,83.03650000)
(122,83.18750000)
(123,83.21550000)
(124,83.29100000)
(125,83.40700000)
(126,84.08500000)
(127,85.71700000)
(128,85.73550000)
(129,85.76850000)
(130,85.80400000)
(131,85.81550000)
(132,85.83950000)
(133,85.99750000)
(134,86.00400000)
(135,86.02200000)
(136,86.03950000)
(137,86.18050000)
(138,86.20400000)
(139,86.24300000)
(140,86.35150000)
(141,86.56300000)
(142,88.22500000)
(143,88.23500000)
(144,88.26350000)
(145,88.29750000)
(146,88.30650000)
(147,88.32300000)
(148,88.45650000)
(149,88.46800000)
(150,88.47750000)
(151,88.49800000)
(152,88.53950000)
(153,88.63450000)
(154,88.67050000)
(155,88.76900000)
(156,88.90250000)
(157,90.31450000)
(158,90.31950000)
(159,90.33800000)
(160,90.36550000)
(161,90.38000000)
(162,90.39250000)
(163,90.49150000)
(164,90.50200000)
(165,90.50950000)
(166,90.52300000)
(167,90.54700000)
(168,90.60900000)
(169,90.62700000)
(170,90.72300000)
(171,90.80150000)
(172,91.92600000)
(173,91.93300000)
(174,91.94900000)
(175,91.97200000)
(176,91.97950000)
(177,91.99500000)
(178,92.00750000)
(179,92.07500000)
(180,92.08000000)
(181,92.08950000)
(182,92.11100000)
(183,92.16950000)
(184,92.18250000)
(185,92.25200000)
(186,92.31550000)
(187,93.24600000)
(188,93.25500000)
(189,93.27000000)
(190,93.28250000)
(191,93.29250000)
(192,93.30150000)
(193,93.31150000)
(194,93.38750000)
(195,93.39150000)
(196,93.39850000)
(197,93.42000000)
(198,93.47250000)
(199,93.48350000)
(200,93.54000000)
(201,93.55700000)
(202,94.42550000)
(203,94.43650000)
(204,94.44050000)
(205,94.45000000)
(206,94.46250000)
(207,94.47150000)
(208,94.48300000)
(209,94.53650000)
(210,94.54200000)
(211,94.55100000)
(212,94.55800000)
(213,94.61050000)
(214,94.62200000)
(215,94.64150000)
(216,94.68450000)
(217,94.78000000)
(218,95.33000000)
(219,95.33350000)
(220,95.34100000)
(221,95.35550000)
(222,95.36350000)
(223,95.37150000)
(224,95.41400000)
(225,95.41850000)
(226,95.42400000)
(227,95.43900000)
(228,95.46450000)
(229,95.49900000)
(230,95.51150000)
(231,95.54550000)
(232,95.60150000)
(233,96.08450000)
(234,96.08650000)
(235,96.09550000)
(236,96.10500000)
(237,96.11000000)
(238,96.11200000)
(239,96.14700000)
(240,96.15200000)
(241,96.15600000)
(242,96.16200000)
(243,96.17350000)
(244,96.20250000)
(245,96.21200000)
(246,96.24900000)
(247,96.28300000)
(248,96.71700000)
(249,96.72050000)
(250,96.72650000)
(251,96.73200000)
(252,96.73550000)
(253,96.74100000)
(254,96.74850000)
(255,96.77900000)
(256,96.78100000)
(257,96.78650000)
(258,96.79500000)
(259,96.82050000)
(260,96.82600000)
(261,96.85750000)
(262,96.88550000)
(263,97.19750000)
(264,97.20100000)
(265,97.20400000)
(266,97.20850000)
(267,97.21450000)
(268,97.21800000)
(269,97.22050000)
(270,97.24650000)
(271,97.24800000)
(272,97.25400000)
(273,97.26050000)
(274,97.28050000)
(275,97.28300000)
(276,97.30600000)
(277,97.31400000)
(278,97.64200000)
(279,97.64550000)
(280,97.64650000)
(281,97.65500000)
(282,97.65900000)
(283,97.66050000)
(284,97.66500000)
(285,97.68800000)
(286,97.69100000)
(287,97.69550000)
(288,97.69950000)
(289,97.72000000)
(290,97.72450000)
(291,97.73550000)
(292,97.76300000)
(293,97.80250000)
(294,98.00050000)
(295,98.00100000)
(296,98.00300000)
(297,98.00650000)
(298,98.00850000)
(299,98.01150000)
(300,98.03600000)
(301,98.03800000)
(302,98.04100000)
(303,98.04550000)
(304,98.05400000)
(305,98.06600000)
(306,98.07450000)
(307,98.08700000)
(308,98.11950000)
(309,98.31200000)
(310,98.31200000)
(311,98.31350000)
(312,98.31700000)
(313,98.32050000)
(314,98.32200000)
(315,98.33650000)
(316,98.33800000)
(317,98.33850000)
(318,98.34300000)
(319,98.35200000)
(320,98.36800000)
(321,98.37200000)
(322,98.39250000)
(323,98.41100000)
(324,98.57800000)
(325,98.58050000)
(326,98.58300000)
(327,98.58600000)
(328,98.58750000)
(329,98.58900000)
(330,98.59050000)
(331,98.60150000)
(332,98.60300000)
(333,98.60550000)
(334,98.60850000)
(335,98.61550000)
(336,98.61800000)
(337,98.62900000)
(338,98.63750000)
(339,98.77950000)
(340,98.78000000)
(341,98.78050000)
(342,98.78300000)
(343,98.78600000)
(344,98.78800000)
(345,98.79000000)
(346,98.80100000)
(347,98.80100000)
(348,98.80450000)
(349,98.80800000)
(350,98.81600000)
(351,98.81750000)
(352,98.83150000)
(353,98.83650000)
(354,98.96400000)
(355,98.96600000)
(356,98.96700000)
(357,98.96750000)
(358,98.96900000)
(359,98.97050000)
(360,98.97100000)
(361,98.98300000)
(362,98.98350000)
(363,98.98400000)
(364,98.98550000)
(365,98.99800000)
(366,98.99900000)
(367,99.00300000)
(368,99.01100000)
(369,99.03250000)
(370,99.11450000)
(371,99.11550000)
(372,99.11550000)
(373,99.11900000)
(374,99.12100000)
(375,99.12600000)
(376,99.13300000)
(377,99.13300000)
(378,99.13500000)
(379,99.13750000)
(380,99.13950000)
(381,99.14250000)
(382,99.14500000)
(383,99.15250000)
(384,99.16000000)
(385,99.23150000)
(386,99.23150000)
(387,99.23350000)
(388,99.23650000)
(389,99.23700000)
(390,99.23750000)
(391,99.24600000)
(392,99.24650000)
(393,99.24650000)
(394,99.24700000)
(395,99.25100000)
(396,99.25750000)
(397,99.25950000)
(398,99.26600000)
(399,99.27500000)
(400,99.33750000)
};
\addlegendentry{別解が真}
% END_EMPIRICAL_CDF
\end{axis}
\end{tikzpicture}
\captionof{figure}{固定した2候補の将来比較。曲線は保存済み各20万系列の経験累積率であり、習得状態の信用区間ではない。}\label{fig:future}
\end{center}

経験中央値は基準が真なら38問、別解が真なら8問であるが、75.000\%分位点は77問／95問となる。
100問以内の到達率は83.518\%／78.915\%、そのモンテカルロ標準誤差（MCSE）は0.083／0.091パーセントポイントであった。
初期に証拠が得られやすい候補でも裾は長く、中央値だけで一般的な必要問題数は決められない。
これは仮定した候補間の比較であり、実際の習得状態を直接観測した検証ではない。

\subsection{保証する対象を、潜在状態から観測予測へ移す}
第$t$問の出題後・解答前の情報を$\F_{t-1}$とし、問題・提示条件・過去の回答・モデルの版を含める。
解答前の正答予測$r_t$と賭け率$\lambda_t$を確定し、
\begin{equation}
 H_{\rm cal}:\ E[Y_t\given\F_{t-1}]=r_t,\qquad
 -\frac{1}{1-r_t}\le\lambda_t\le\frac{1}{r_t}\quad(0<r_t<1)
 \label{eq:calnull}
\end{equation}
を考える。これは予測値ごとの平均が合うという較正より強い、履歴と出題条件を含む条件付き平均の仮説である。
事前に定めた$0<\alpha<1$に対し、初期値$E_0=1$の過程は\cite{betting}
\begin{equation}
 E_n=\prod_{t=1}^n\{1+\lambda_t(Y_t-r_t)\},\qquad
 P_{H_{\rm cal}}\!\left(\exists n:E_n\ge1/\alpha\right)\le\alpha
 \label{eq:caltest}
\end{equation}
が成り立つ。因子の条件付き期待値が1、非負性が式\eqref{eq:calnull}で保たれるためである。
式\eqref{eq:caltest}と式\eqref{eq:ville}は、帰無仮説の下での非負マルチンゲールにVilleの不等式を適用する点で共通する。
賽では通常賽のみという仮説、学習では式\eqref{eq:calnull}の条件付き平均を帰無仮説とする。
比較予測$s_t$を解答前に定め、$\lambda_t=(s_t-r_t)/\{r_t(1-r_t)\}$とすれば、各因子はBernoulli予測の尤度比となり、前節の将来比較とつながる。
適応的出題でも、選ばれた問題を条件に予測し、上記仮説と解答前の確定を満たすことが必要である。

検算用の一定予測$r_t=0.800$、比較予測$s_t=0.600$、境界100では、正答で$3/4$倍、誤答で2倍となる。
整数再帰による100問以内の誤警報率は0.712\%、真の正答率0.600での検知率は93.522\%である。
この性能値は時変BKT一般の値ではない。期間無制限の誤警報上界1.000\%と、有限期間の性能も区別する。
\textbf{無警報は予測の正しさの証明ではなく、潜在習得状態や誤卒業率への直接の保証でもない。}

\newpage
\section{不確かさの評価と、実社会への応用}
\subsection{推定の幅は、発生源と評価方法を分けて扱う}
状態や未知パラメータの不確かさは、その推定確率を求める計算誤差と区別する。
未知パラメータ$\theta$を含む習得確率は、例えば
\begin{equation}
 \bar p_t=\int P(L_t=1\given D,\theta)\,p(\theta\given D)\,d\theta
\end{equation}
と平均化できる。状態・パラメータの不確かさ、計算の精度、モデル外のずれは別々に評価する。
今回の固定パラメータBKTは厳密再帰、将来到達率は独立系列のモンテカルロで評価し、MCMCは用いていない。

\begin{center}
\captionof{table}{不確かさ・誤差の評価を分ける}\label{tab:uncertainty}
\small
\begin{tabular}{@{}>{\raggedright\arraybackslash}p{34mm}>{\raggedright\arraybackslash}p{54mm}>{\raggedright\arraybackslash}p{64mm}@{}}
\toprule 対象 & 例 & 評価の方法\\
\midrule
モデル内の不確かさ & 正誤の変動、未知の習得状態、有限データからのパラメータ推定 & 事後分布、事後予測、将来課題での較正、事前分布への感度\\
近似計算の誤差 & 有限回の模擬実験、数値積分、MCMCを使う場合の有限長 & MCSE、厳密解・独立実装との照合。MCMCなら収束・有効標本数も確認\\
モデル・測定等のずれ & 難易度、誤採点、スキル付与、忘却、欠測、出題方針 & 共通項目、条件別・時系列の外部検証、モデルへの感度分析\\
\bottomrule
\end{tabular}
\end{center}
既知真値の模擬実験で手順を確かめ、学習者・項目・時系列を分けて将来予測を検証する。
自信度・解答時間は測定の意味と依存関係を定めて追加する。
感度比較を信用区間と呼ばず、MCSEを状態確率の誤差へ流用しない。

\subsection{学習支援・LLM品質管理と、因果効果の検証}
学習支援基盤では観測と推定結果を分離し、解答前予測、出題条件、問題・モデル等の版と結果を保存する。
支援判断は共通の将来課題、転移・定着、誤判断の損失も含めて評価する。
LLM生成問題では正確さ・正解の一意性・スキル適合・難易度を検証し、誤答をすべて未習得へ帰さない。
同じ生成元やLLMによる判定の重複計上も避ける。

解説$A$を変える因果効果は、将来の共通課題$Y$を用いて
\begin{equation}
 \Delta=E[Y(a)]-E[Y(a')]
\end{equation}
と定め、無作為割付、または一貫性・交換可能性・正値性等の識別条件を検討する\cite{causal}。
事前分布だけで未観測交絡は解消しない。学習者データ、LLM品質管理、介入効果は今後の実証課題である。

\section{結論}
地下チンチロのベイズ推定を土台に、外的要因・共有状態・学習遷移を取り込む推論を検討した。
賽の仮説判定では、指定した事前確率と閾値の下で期間無制限の誤警報上界$1/891$（約0.112\%）を得た。
理解度推定では、同じ8回答への尤度がほぼ等しくても習得確率は0.984と0.369に分かれた。
潜在状態の推定と観測予測の保証を区別し、後者を非負マルチンゲールによって監視する方法を示した。
これらは、見えない状態や指標を扱う他分野への応用につながる。
実用性は、対象ごとの観測設計と実データで検証する。

\vspace{1mm}
\noindent{\footnotesize\color{muted}\textbf{資料}\enspace
数値表示は小数点以下3桁へ四捨五入し、図の内部データは丸めず保持した。}

\begin{thebibliography}{6}
\footnotesize\setlength{\itemsep}{0pt}\setlength{\parskip}{0pt}
\bibitem{kaiji} 福本伸行『賭博破戒録カイジ』、講談社、地下チンチロ編。物語の着想の出典。
\bibitem{howard} S. R. Howard et al., ``Time-uniform, nonparametric, nonasymptotic confidence sequences,'' \textit{Ann. Statist.} 49, 1055--1080 (2021). \href{https://doi.org/10.1214/20-AOS1991}{doi:10.1214/20-AOS1991}.
\bibitem{corbett} A. T. Corbett and J. R. Anderson, ``Knowledge tracing: Modeling the acquisition of procedural knowledge,'' \textit{UMUAI} 4, 253--278 (1995). \href{https://doi.org/10.1007/BF01099821}{doi:10.1007/BF01099821}.
\bibitem{beck} J. E. Beck and K.-m. Chang, ``Identifiability: A Fundamental Problem of Student Modeling,'' \textit{UM 2007}, 137--146 (2007).
\bibitem{betting} I. Waudby-Smith and A. Ramdas, ``Estimating means of bounded random variables by betting,'' \textit{J. R. Stat. Soc. B} 86(1), 1--27 (2024), §4. \href{https://doi.org/10.1093/jrsssb/qkad009}{doi:10.1093/jrsssb/qkad009}.
\bibitem{causal} M. A. Hern\'an and J. M. Robins, \textit{Causal Inference: What If} (2020). \url{https://miguelhernan.org/whatifbook}.
\end{thebibliography}
\end{document}
